% SOURCE_AUTHORITY: sga5_fr_workpass.tex, lines 7446--7464 (LF-normalized unit).
% SOURCE_SHA256: 791F4EFFC5E02832D5D77ED03518C8156D6F07E4C8238B03545DB93D883FBB28
\begin{statement}{Proposição 4.2.5}
Caracterização dos sistemas projetivos estritos AR-\(J\)-ádicos por meio dos graduados associados. Seja \(X\) um sistema projetivo estrito indexado por \(\mathbb N\), com valores em \(\mathcal C\). Supõe-se que, para todo inteiro \(n\), tem-se
\[
J^{n+1}X_n=0.
\]
As seguintes afirmações são equivalentes:
\begin{enumerate}[label=\textup{(\roman*)}]
\item \(X\) é AR-\(J\)-ádico.
\item Existe um inteiro \(r>0\) tal que, para todo inteiro \(n\), verifica-se a inclusão que se encontra abaixo:
\[
F^n(X_n)\subset J^{\,n-r}X_n.
\]
\item Existe um inteiro \(r\) tal que, para todo inteiro \(m\), o morfismo canônico que aparece abaixo é um isomorfismo sempre que \(n\ge m+r\):
\[
\gr^n(J^mX)\longrightarrow \gr^n(X).
\]
\end{enumerate}
No enunciado de (ii), convém pôr \(J^{\,n-r}X_n=X_n\) quando \(n<r\).
\end{statement}
